\section{Method and Theory}
\label{sec:method}

\subsection{Black--Rangarajan dual and the GNC-QPEP forward pass}

\begin{lemma}[exact BR dual for GM]\label{lem:brdual}
For every $\mu\ge1$, $c>0$, and $e\ge0$,
\begin{equation}
\begin{aligned}
\rho_\mu(e)&=\min_{w\in[0,1]}\big\{we+\psi_\mu(w)\big\},\\
w^*(e)&=\rho_\mu'(e)=\Big(\frac{\mu c^2}{\mu c^2+e}\Big)^2,
\end{aligned}
\end{equation}
and consequently
$\min_{w\in[0,1]^N}\sum_i\{w_ie_i(z)+\psi_\mu(w_i)\}=\sum_i\rho_\mu(e_i(z))$ for every $z$.
The scalar objective is strictly convex on $(0,1]$ with stationary point $w^*$, and the box
constraint is a Cartesian product; the complete proof is in Appendix~\ref{app:proofs}.
\end{lemma}

\begin{theorem}[GNC-QPEP forward]\label{thm:gnc}
Under (A1), (A2), (A5), fix $\mu\ge1$ and alternate
\begin{equation}
w_i^{(k)}=\rho_\mu'\big(e_i(z^{(k)})\big),\qquad
z^{(k+1)}=\argmin_{z\in\mathcal{M}}\sum_iw_i^{(k)}e_i(z).
\end{equation}
Then:
\begin{enumerate}
\item[(i)] The $w$-step is exact and the $z$-step is a weighted QPEP. Eliminating $\mathbf{t}$ gives
  $\mathbf{t}(w,\mathbf{q})=G(w)^{-1}\sum_iw_iW_i(\mathbf{y}_i-R(\mathbf{q})\mathbf{p}_i)$ with
  $G(w)=\sum_iw_iW_i$; the reduced cost is a quartic polynomial in $\mathbf{q}$, the KKT system has
  four (or seven) variables, and $N$ enters only through $O(N)$ accumulation into fixed-size
  matrices.
\item[(ii)] $F_\mu(z^{(k)})$ is nonincreasing; if each weighted minimizer is unique, the iterates
  converge to a fixed point $(z_\mu^*,w_\mu^*)$ with $w_{\mu,i}^*=\rho_\mu'(e_i(z_\mu^*))$ and
  $z_\mu^*\in\argmin_z\sum_iw_{\mu,i}^*e_i(z)$.
\item[(iii)] As $\mu\downarrow1$, every accumulation point satisfies $w_i^*=\rho'(e_i(z^*))$ and
  $\Phi(z^*)=0$, i.e.\ $z^*$ is a stationary point of $F_{\mathrm{rob}}$.
\item[(iv)] The global property belongs to the inner weighted QPEP only: the outer GNC is a
  nonconvex block-coordinate method that can stop at a saddle, and a coherent majority of outliers
  can make the global minimizer of $F_{\mathrm{rob}}$ itself wrong
  (Section~\ref{sec:discussion}).
\end{enumerate}
Statement (ii) is the standard monotone block-descent argument; (iii) follows by continuity of
$\rho_\mu'$ and the first-order condition of the weighted problem. In the 3D--3D experiments the
inner solver is the closed-form scalar-weight $4\times4$ specialization of QPEP; a
$4\times4$ Davenport/Horn variant is benchmarked in Appendix~\ref{app:systems}.
\end{theorem}

\begin{algorithm}[t]
\caption{Forward GNC-QPEP (one robust solve)}
\label{alg:forward}
\footnotesize
\begin{algorithmic}[1]
\Require $\{(\mathbf{p}_i,\mathbf{y}_i,W_i)\}_{i=1}^N$, scale $c$, schedule $\gamma$
\State initialize $\mathbf{q},\mathbf{t}$; $\mu\gets3\max_ie_i/c^2$
\While{$\mu\ge1$}
  \For{$k=1$ to $K$}
    \State $e_i\gets\|R(\mathbf{q})\mathbf{p}_i+\mathbf{t}-\mathbf{y}_i\|_{W_i}^2$ for all $i$
    \State $w_i\gets\big(\mu c^2/(\mu c^2+e_i)\big)^2$ for all $i$
    \State $(\mathbf{q},\mathbf{t})\gets\textsc{QPEP-Global}\big(\sum_iw_ie_i\big)$
  \EndFor
  \State $\mu\gets\max(1,\mu/\gamma)$
\EndWhile
\State \Return $\mathbf{q},\mathbf{t},w,e$
\end{algorithmic}
\end{algorithm}

The annealing schedule is $\mu_0\ge3\max_ie_i(z^{(0)})/c^2$,
$\mu_{k+1}=\max(1,\mu_k/\gamma)$ with $\gamma\in[1.2,2.0]$, and a fixed or adaptive number of inner
alternations. The choice of $\mu_0$ makes each residual's \emph{radial} second derivative positive;
it does not make the full cost convex in $z$.

\subsection{Fixed points, curvature correction, and the exact Riemannian Hessian}

\begin{lemma}[fixed point $\Rightarrow$ stationarity]\label{lem:fixedpoint}
If $z^*$ globally minimizes $\sum_iw_i^*e_i$ with $w_i^*=\rho'(e_i(z^*))$, then $(z^*,w^*)$ is a
GNC fixed point and $\Phi(z^*)=0$. Conversely, $\Phi(z^*)=0$ plus the same weighted-global-minimum
property gives a fixed point; stationarity alone is not enough.
\end{lemma}

\begin{lemma}[curvature correction $=$ Schur complement]\label{lem:schur}
At a GNC fixed point, let $H_{\mathrm{IRLS}}=\sum_i\rho'(e_i^*)H_i$,
$\Delta H_{\mathrm{redesc}}=\sum_i\rho''(e_i^*)\,g_ig_i^\top\preceq0$, $D=[g_1,\dots,g_N]$ and
$\Psi''=\diag(\psi''(w_i^*))$. The Hessian of the joint BR objective is
\begin{equation}
\begin{aligned}
H_J&=\begin{bmatrix}H_{\mathrm{IRLS}}&D\\D^\top&\Psi''\end{bmatrix},\\
H_{\mathrm{rob}}&=H_{\mathrm{IRLS}}-D\Psi''^{-1}D^\top
=H_{\mathrm{IRLS}}+\Delta H_{\mathrm{redesc}}.
\end{aligned}
\end{equation}
Since $\psi''(w_i^*)=-1/\rho''(e_i^*)>0$, the redescending correction is exactly the eliminated
weight block, and if $H_{\mathrm{IRLS}}\succeq0$ then $H_J\succ0\iff H_{\mathrm{rob}}\succ0$.
Freezing the weights deletes this Schur complement; Section~\ref{sec:gradient} quantifies the
resulting gradient error.
\end{lemma}

\begin{lemma}[exact Riemannian Hessian on $\mathbb{S}^3\times\mathbb{R}^3$]\label{lem:hessian}
Let $f\in C^2(\mathcal{M})$ with Euclidean gradient $G$ and Hessian $H_E$. Then
$\operatorname{Hess}f=\Pi\big(H_E-(\mathbf{q}^\top G_q)\diag(I_4,0_3)\big)\Pi$, and in an
orthonormal tangent basis $B$,
\begin{equation}
H_f=B^\top H_E B-(\mathbf{q}^\top G_q)\diag(I_3,0_3).
\end{equation}
For $F_{\mathrm{rob}}$,
\begin{equation}
\boxed{
\begin{aligned}
H_{\mathrm{rob}}&=B^\top\Big[\sum_i\big(\rho'(e_i)H_i^E+\rho''(e_i)g_i^E(g_i^E)^\top\big)\Big]B\\
&\quad-\lambda_{\mathrm{rob}}\diag(I_3,0_3),\\
\lambda_{\mathrm{rob}}&=\mathbf{q}^\top G_q=\sum_i\rho'(e_i)\lambda_i .
\end{aligned}}
\end{equation}
If $H_i$ and $g_i$ denote Riemannian quantities, the equivalent form is
$H_{\mathrm{rob}}=\sum_i\rho'(e_i)H_i+\sum_i\rho''(e_i)g_ig_i^\top$. The term
$-\lambda_{\mathrm{rob}}\diag(I_3,0_3)$ is the multiplier of the unit-quaternion constraint and
cannot be absorbed into a rescaling: for 3D--3D,
\begin{equation}
\lambda_i=\mathbf{q}^\top g_{i,q}^E=4\mathbf{u}_i^\top W_i\mathbf{r}_i
=4\big(e_i-(\mathbf{t}-\mathbf{y}_i)^\top W_i\mathbf{r}_i\big)\neq4e_i
\end{equation}
whenever $(\mathbf{t}-\mathbf{y}_i)^\top W_i\mathbf{r}_i\neq0$. On the sphere,
$f(\mathbf{q})=\frac12\|\mathbf{q}\|^2$ is constant and has zero Riemannian Hessian, whereas the
Euclidean formula gives $I_3$ in tangent coordinates. In a six-dimensional 3D--3D example with
$c=1.0$, the uncorrected formula differs from a geodesic finite-difference Hessian by $25.9\%$ in
relative Frobenius norm, while the corrected formula agrees to $9.9\times10^{-8}$; the multiplier
is $\lambda_{\mathrm{rob}}=1.033$.
\end{lemma}

\begin{lemma}[3D--3D weights]\label{lem:weights}
With the notation above, $g_{i,q}^E=2J_i^\top W_i\mathbf{r}_i$, $g_{i,t}^E=2W_i\mathbf{r}_i$, and
$\lambda_i=4\mathbf{u}_i^\top W_i\mathbf{r}_i$. The total multiplier is
$\lambda_{\mathrm{rob}}=4\sum_i\rho'(e_i)\mathbf{u}_i^\top W_i\mathbf{r}_i$, not
$4\sum_i\rho'(e_i)e_i$.
\end{lemma}

\subsection{Exact implicit differentiation}

\begin{theorem}[exact fixed-point IFT]\label{thm:ift}
Under (A1)--(A4), suppose $\Phi(z^*,\theta_0)=0$ and
$H_{\mathrm{rob}}(z^*,\theta_0)$ is invertible. Then there is a neighborhood $U$ of $\theta_0$ and
a unique $C^1$ branch $\bar z:U\to\mathcal{M}$ with $\bar z(\theta_0)=z^*$ and
$\Phi(\bar z(\theta),\theta)=0$, and
\begin{equation}
\frac{\mathrm{d}\bar z}{\mathrm{d}\theta}=-H_{\mathrm{rob}}^{-1}\frac{\partial\Phi}{\partial\theta}.
\end{equation}
If $H_{\mathrm{rob}}\succ0$ the branch is a strict local minimum and remains so for small
perturbations; a negative eigenvalue makes $z^*$ a saddle; a singular $H_{\mathrm{rob}}$ makes the
implicit-function theorem fail, and the branch can bifurcate. The derivative is that of \emph{the
selected stationary branch}; if the GNC iteration switches branches, the derivative refers to the
branch it selects. Invertibility is the IFT condition, while $\lambda_{\min}>0$ is the stronger
condition for a strict local minimum and branch persistence.
\end{theorem}

\begin{theorem}[right-hand side decomposition]\label{thm:rhs}
For a parameter $\theta$,
\begin{equation}
\frac{\partial\Phi}{\partial\theta}
=\underbrace{\sum_i\rho'(e_i)\frac{\partial g_i}{\partial\theta}}_{\text{frozen-weight term}}
+\underbrace{\sum_i\rho''(e_i)\,g_i\Big(\frac{\partial e_i}{\partial\theta}\Big)^\top}
_{\text{implicit weight sensitivity}}.
\end{equation}
For $\theta=\mathbf{y}_i$ in 3D--3D,
\begin{equation}
\begin{aligned}
\frac{\partial\Phi_t}{\partial\mathbf{y}_i}&=-2\big(\rho'(e_i)W_i
+2\rho''(e_i)W_i\mathbf{r}_i\mathbf{r}_i^\top W_i\big),\\
\frac{\partial\Phi_q}{\partial\mathbf{y}_i}&=-2J_i^\top W_i
\big(\rho'(e_i)I_3+2\rho''(e_i)\mathbf{r}_i\mathbf{r}_i^\top W_i\big).
\end{aligned}
\end{equation}
The frozen-weight approximation keeps only the first term and sets
$\partial\Phi_t^{\mathrm{naive}}/\partial\mathbf{y}_i=-2\rho'(e_i)W_i$.
\end{theorem}

\subsection{When the frozen-weight gradient has the wrong sign}
\label{sec:signflip}

With $W_i=I_3$, the parallel eigenvalue of $\partial\Phi_t/\partial\mathbf{y}_i$ is
$-\lambda_\parallel(e_i)$, where
\begin{equation}
\lambda_\parallel(e)=2(\rho'(e)+2e\rho''(e))=\frac{2c^4(c^2-3e)}{(c^2+e)^3}.
\end{equation}
Thus $\lambda_\parallel(0)=2$, $\lambda_\parallel(c^2/3)=0$, $\lambda_\parallel(c^2)=-1/2$, and
$\lambda_\parallel(e)\sim-6c^4/e^2$ as $e\to\infty$. The exact and naive parallel eigenvalues differ
by $-4e\rho''=8c^4e/(c^2+e)^3>0$, so for $e>c^2/3$ they have strictly opposite signs; the most
influential residuals are the \emph{boundary} ones near $e=c^2$, not the most distant gross
outliers.

\begin{proposition}[1D sign flip]\label{prop:signflip}
Let $F(t)=\sum_i\rho((t-y_i)^2)$ with GM. At a stationary point $t^*$ with
$H_{\mathrm{rob}}=\sum_i\lambda_\parallel(e_i)>0$,
\begin{equation}
\frac{\mathrm{d}t^*}{\mathrm{d}y_j}=\frac{\lambda_\parallel(e_j)}{H_{\mathrm{rob}}},
\qquad
\Big(\frac{\mathrm{d}t^*}{\mathrm{d}y_j}\Big)_{\mathrm{naive}}
=\frac{2\rho'(e_j)}{\sum_i2\rho'(e_i)}>0.
\end{equation}
If $e_j>c^2/3$, then $\lambda_\parallel(e_j)<0$, so the exact derivative is strictly negative while
the naive derivative is strictly positive. If $H_{\mathrm{rob}}=0$ the implicit-function theorem
fails; if $H_{\mathrm{rob}}<0$ the point is not a local minimum.
\end{proposition}

\begin{proposition}[existence]\label{prop:existence}
For any $c>0$ and $y_2-y_1$ sufficiently large, the two-observation cost has a strict local minimum
$t^*$ with $e_1<c^2/3<e_2$ and $H_{\mathrm{rob}}>0$; hence the sign flip of
Proposition~\ref{prop:signflip} occurs. The proof uses the signs of $\Phi$ at $0$ and at $Kc^4/d^3$,
and the limit $H_{\mathrm{rob}}\to\lambda_\parallel(0)=2>0$.
\end{proposition}

In 6D the sign of the pose sensitivity is \emph{not} determined by $\lambda_\parallel$ alone:
$\partial z^*/\partial\mathbf{y}_i=-H_{\mathrm{rob}}^{-1}\partial\Phi/\partial\mathbf{y}_i$, and the
inverse Hessian mixes directions. Proposition~\ref{prop:signflip} is the rigorous statement; the 6D
counterpart is numerical, with a cosine of $-0.079$, and a 6D analogue would require a joint
spectral condition on $H_{\mathrm{rob}}$ and $\partial\Phi/\partial\mathbf{y}_i$.

\begin{proposition}[per-observation $O(\|\mathbf{r}\|^{-4})$ bound]\label{prop:bound}
Under (A1), (A5), assume along a path $\|\mathbf{y}_i\|\to\infty$ that
$z^*(\mathbf{y}_i)$ stays in a compact set and $H_{\mathrm{rob}}$ is uniformly positive definite
there with $\lambda_{\min}\ge\mu_0>0$. Then
\begin{equation}
\Big\|\frac{\partial\Phi}{\partial\mathbf{y}_i}\Big\|_F\le C\frac{c^4}{(c^2+e_i)^2}
\le C\frac{c^4}{\|\mathbf{r}_i\|^4},
\end{equation}
\begin{equation}
\Big\|\frac{\partial z^*}{\partial\mathbf{y}_i}\Big\|_F\le\frac{C}{\mu_0}
\frac{c^4}{\|\mathbf{r}_i\|^4}.
\end{equation}
The forward influence is only $O(\|\mathbf{r}_i\|^{-3})$:
$\|\rho'(e_i)g_i\|\le2\overline wc^4/\|\mathbf{r}_i\|^3$. The bound is per observation; it does not
prevent $N$ outliers from accumulating, and it does not constrain a coherent outlier cluster that
becomes the global minimum. Gradient decay through a single observation is therefore distinct from
forward robustness of the estimator.
\end{proposition}

\subsection{Algorithm, complexity, and the second-order certificate}

The backward pass (Algorithm~\ref{alg:backward}, Appendix~\ref{app:backward}) stores $z^*,w^*$ and
recomputes $g_i^E,H_i^E$; it solves one $6\times6$ system per parameter block and never unrolls the
GNC iterations. Memory is $O(N)$ for weights and residuals and $O(1)$ in the number of GNC
iterations; the complexity is $O(N)$ per forward iteration plus a fixed-size algebraic solve, and
$O(N)$ for the backward assembly plus $O(6^3)$. The \textbf{second-order certificate} is the
spectrum of $H_{\mathrm{rob}}$: $H_{\mathrm{rob}}\succ0$ makes the fixed point a strict local
minimum of the joint BR objective, a negative eigenvalue makes it a saddle, and a zero crossing
makes the branch bifurcate and the gradient norm blow up. This is a local stability diagnostic: a
wrong, outlier-consistent pose can also be a strict local minimum with $H_{\mathrm{rob}}\succ0$.
